% !TeX program = lualatex
% Compile twice: lualatex 110.tex
% All references and prose are embedded; no BibTeX database or external inputs.
\documentclass[11pt,a4paper]{article}
\usepackage[margin=24mm,headheight=14pt]{geometry}
\usepackage{fontspec}
\setmainfont{Latin Modern Roman}
\setsansfont{Latin Modern Sans}
\setmonofont{Latin Modern Mono}
\usepackage{amsmath,amssymb,mathtools}
\usepackage{unicode-math}
\setmathfont{Latin Modern Math}
\usepackage{microtype}
\usepackage{enumitem,xurl,needspace}
\usepackage[dvipsnames]{xcolor}
\usepackage{hyperref}
\hypersetup{unicode=true,colorlinks=true,linkcolor=MidnightBlue,urlcolor=MidnightBlue,
 pdftitle={Research attempt: Problem 110},pdfauthor={Research notebook},bookmarksnumbered=true}
\usepackage{bookmark}
\usepackage{tocloft}
\setlength{\cftsecnumwidth}{3em}
\cftsetpnumwidth{2.5em}
\cftsetrmarg{4em}
\usepackage{titlesec}
\titleformat{\section}{\Large\bfseries\raggedright}{\thesection}{0.7em}{}
\usepackage{fancyhdr}
\pagestyle{fancy}\fancyhf{}
\fancyhead[L]{\small\sffamily Open Problems Math | Research notebook}
\fancyhead[R]{\small Problem 110 | 27 September 2026}
\fancyfoot[C]{\thepage}
\setlength{\parindent}{0pt}
\setlength{\parskip}{3pt plus 1pt}
\setlength{\emergencystretch}{3em}
\setcounter{tocdepth}{1}
\setcounter{secnumdepth}{1}
\setlist[itemize]{leftmargin=3em,itemsep=5pt,topsep=4pt,parsep=0pt}
\allowdisplaybreaks
\newcommand{\N}{\mathbb{N}}
\newcommand{\Z}{\mathbb{Z}}
\newcommand{\Q}{\mathbb{Q}}
\newcommand{\R}{\mathbb{R}}
\newcommand{\C}{\mathbb{C}}
\newcommand{\F}{\mathbb{F}}
\newcommand{\A}{\mathbb{A}}
\newcommand{\PP}{\mathbb P}
\newcommand{\E}{\mathbb{E}}
\newcommand{\eps}{\varepsilon}
\newcommand{\dd}{\,\mathrm d}
\newcommand{\sref}[2]{\hyperlink{source:#1:#2}{[S#2]}}
\newcommand{\eref}[2]{\hyperlink{extra:#1:#2}{[E#2]}}
% Turn the next switch off for a shorter reading copy. The quotations preserve
% every exactTarget field, including qualifications omitted from a short summary.
\newif\ifshowcatalogue
\showcataloguetrue
\newcommand{\cataloguescope}[1]{\ifshowcatalogue
 \paragraph{Exact scope in the supplied catalogue.}
 \begin{quote}\small #1\end{quote}\fi}
\usepackage{etoolbox}
\AtBeginEnvironment{itemize}{\small}
\numberwithin{equation}{section}
% Standalone export. Original section 110 preserved verbatim outside marked additions.
\begin{document}
\setcounter{section}{109}
% ================================================================
% problemId: problem.coding-theory-and-combinatorics-problem-cc-8-exact-kissing-numbers-in-higher-dimensions-n-4
% source releaseRank: 110; source releaseStatus: open_with_solved_subcases
% exactTarget SHA-256: e4ae933774fe5898e8cbdfecae75c499d842781b86e3ed72599c2bd82b32426a
\clearpage
\section{\texorpdfstring{Exact Kissing Numbers in Higher Dimensions}{Exact Kissing Numbers in Higher Dimensions}}\label{110}
\subsection{Definitions and mathematical statement}
After scaling the centers of touching unit balls to the unit sphere, the kissing number has the equivalent finite-dimensional definition
\[
 \tau_n=\max\bigl\{|X|:X\subset S^{n-1},\ \langle x,y\rangle\le\tfrac12
                                  \text{ for distinct }x,y\in X\bigr\}.
\]
The inner-product condition is equivalent to angular separation at least $\pi/3$. The problem asks for the exact integer $\tau_n$ for every $n>4$, not only asymptotic bounds or the best lattice configuration. Centers may be arranged arbitrarily, subject to nonoverlap.
\par\smallskip\textit{Formulation and concept references:} \sref{110}{1}.
\cataloguescope{For every integer n > 4, determine the exact kissing number tau\_n: the maximum number of pairwise nonoverlapping unit balls in Euclidean R\textasciicircum{}n that can touch a central unit ball of the same radius.}
\subsection{Short English statement}
Exactly how many equal balls can simultaneously touch one equal central ball in each Euclidean dimension greater than four?
% BEGIN RESEARCH ADDITION 110
\subsection{Research attempt}

\paragraph{Current frontier.}
The target is the entire function $n\mapsto\tau_n$ for $n>4$, not just
an improvement in one dimension. The source-defined spherical-code model
is unambiguous. Cohn--Rajagopal, Section~2 and Theorem~2.1, give four
nonisometric $40$-point configurations in dimension five, while retaining
$40\leq\tau_5\leq44$. Thus even a complete classification of perturbations
of $D_5$ would not classify all plausible optimizers. \sref{110}{2}

Recent lower-bound work remains genuinely constructive rather than an
exact determination in all dimensions: the September 2026 version of
Cohn--Li improves dimensions $17$--$21$ and acknowledges Ho's subsequent
improvement in dimension $19$. A September 2026 preprint of Takhanov--Yun
reports additional higher-dimensional constructions. Their newest claims
are not independently certified by this notebook and are not used in any
proof below. The local/global distinction is already explicit in the
rigidity theory of Cohn--Jiao--Kumar--Torquato. \eref{110}{1}\,\eref{110}{2}\,\eref{110}{3}

\paragraph{Attack route.}
Three possibilities were tested: add points into holes of a root
configuration; deform that configuration to create a new hole; or produce
a global positive-definite-polynomial upper certificate. The first needs
a precise covering-radius calculation, the second a nontrivial feasible
infinitesimal displacement, and the third a bound below the next integer.
We compute all three diagnostics. They yield exact restricted statements
and a particularly concrete counterexample to the inference that
``saturated and rigid'' means ``maximum cardinality.''

\paragraph{Attempt: exact holes and a dimension threshold.}
For $n\geq4$ let
\[
 X_n=\{(\sigma e_i+\tau e_j)/\sqrt2:
           1\leq i<j\leq n,\ \sigma,\tau\in\{-1,1\}\}.
\]
These $2n(n-1)$ unit vectors have pairwise inner products at most $1/2$.
For a unit vector $v$, let $a_1\geq\cdots\geq a_n\geq0$ be its sorted
absolute coordinates. Then
\begin{equation}\label{110:support}
 \max_{x\in X_n}\langle v,x\rangle=(a_1+a_2)/\sqrt2.
\end{equation}
Writing $c=a_1+a_2$, one has
\[
 1=\sum_i a_i^2\leq a_1^2+(n-1)a_2^2
             =a_1^2+(n-1)(c-a_1)^2.
\]
The last expression is convex for $c/2\leq a_1\leq c$, so it is at most
$\max(c^2,nc^2/4)=nc^2/4$ when $n\geq4$. Equality is attained by
$a_i=1/\sqrt n$. Therefore
\begin{equation}\label{110:cover}
 \min_{v\in S^{n-1}}\max_{x\in X_n}\langle v,x\rangle
       =\sqrt{2/n}\qquad(n\geq4).
\end{equation}
For $n>4$, equality forces every absolute coordinate to equal
$1/\sqrt n$: the maximizing endpoint is $a_1=c/2$, and equality must
also hold in $\sum a_i^2\leq a_1^2+(n-1)a_2^2$.

Consequently $X_n$ admits no additional point for $4\leq n\leq7$.
In particular, the exact positive obstruction margin for $D_5$ is
$\sqrt{2/5}-1/2$. This does not merely reflect failure of a numerical
search. Conversely, at $n=8$ the possible extra points are precisely
\[
 v=(\sigma_1,\ldots,\sigma_8)/(2\sqrt2),
 \qquad \sigma_i\in\{-1,1\}.
\]
Two such points have inner product $1-d_H(\sigma,\tau)/4$.
Choose the $128$ sign patterns with an even number of minus signs.
Their mutual Hamming distances are at least two, so adding them to the
$112$ roots gives a valid $240$-point configuration. No larger subset of
these $256$ hole points is possible: pairing each sign pattern with the
one obtained by flipping its first sign partitions them into $128$
forbidden distance-one pairs. This proves optimality only \emph{among
extensions of the fixed $D_8$ configuration} at this stage.

\paragraph{Attempt: an exact infinitesimal certificate.}
Use unnormalized integer roots $r_i=\sqrt2x_i$. A first-order feasible
velocity $v_i$ obeys
\[
 r_i\cdot v_i=0,\qquad
 r_i\cdot v_j+r_j\cdot v_i\leq0
             \quad\text{when }r_i\cdot r_j=1.
\]
For each root,
\[
 \sum_{j:r_i\cdot r_j=1}r_j=2(n-2)r_i.
\]
Summing the contact inequalities thus gives exactly zero, by the tangent
conditions. Since every summand is nonpositive, each contact inequality
is an equality. This is a positive equilibrium stress, proved just by
counting the signed roots that share one coordinate.

The accompanying script builds the integer matrix $M_n$ of these tangent
and contact equalities. Exact Gaussian elimination modulo the prime
$1000003$ gives
\[
 \begin{array}{c|r|r|r|r}
 n&|X_n|&\text{contacts}&\text{variables}&\operatorname{rank}_{\F_{1000003}}M_n\\\hline
 4&24&96&96&90\\
 5&40&240&200&190\\
 6&60&480&360&345
 \end{array}
\]
A nonzero minor modulo a prime is a nonzero integer minor, so these are
rigorous lower bounds for real rank. On the other hand the
$n(n-1)/2$ independent rotations $v_i=Ar_i$, $A^T=-A$, lie in the
kernel. Hence the displayed ranks are exactly the real ranks and all
feasible infinitesimal motions are rotations. This is an exact finite
certificate, not a tolerance-based singular-value calculation. The
conclusion reproduces a known result, not a new rigidity theorem;
Corollary~3.4 of \eref{110}{1} proves it for all $D_n$, $n\geq4$.

\paragraph{Attempt: why the same polynomial closes dimension eight only.}
The diagnostic polynomial
\[
 f(t)=(t+1)(t+\tfrac12)^2t^2(t-\tfrac12)
\]
is nonpositive on $[-1,1/2]$. Let $P_k^{(n)}(1)=1$ be the normalized
Gegenbauer polynomial for the sphere $S^{n-1}$. The addition formula for
spherical harmonics gives
$\sum_{i,j}P_k^{(n)}(\langle x_i,x_j\rangle)\geq0$:
up to a positive normalization it is a sum, over an orthonormal harmonic
basis $Y_{k,l}$, of $|\sum_iY_{k,l}(x_i)|^2$.
Thus, whenever $f=\sum f_kP_k^{(n)}$ with $f_0>0$ and $f_k\geq0$,
\[
 N^2f_0\leq\sum_{i,j}f(\langle x_i,x_j\rangle)\leq Nf(1),
 \qquad N\leq f(1)/f_0.
\]
Exact polynomial expansion in dimension eight gives, in increasing degree,
\[
 (f_0,\ldots,f_6)=
 \left(\frac3{320},\frac3{40},\frac{15}{64},\frac{39}{80},
       \frac{399}{640},\frac9{16},\frac{33}{128}\right).
\]
Since $f(1)=9/4$, this proves $N\leq240$. Together with the explicit
construction it recovers the known equality $\tau_8=240$.
The polynomial identity is independently checked by exact rational
arithmetic in the script.

In dimension five the same polynomial still has positive Gegenbauer
coefficients but $f_0=37/840$, so it gives only
$N\leq1890/37$, hence $N\leq51$. It is worse than the known upper bound.
The attractive contact roots of $D_5$ therefore do \emph{not} manufacture
an optimal dual certificate. A genuinely useful global certificate in
that dimension must exclude cardinality $41$; a bound strictly below $41$
would suffice, but none is produced here.

\paragraph{Stress test: an explicit failure of the local-to-global claim.}
The same calculations show that $D_6$ is saturated and infinitesimally
rigid, yet its $60$ points are not maximum. Here is a direct $72$-point
comparison, without relying on an optimization table. In $\R^6$, use
unnormalized squared norm two and take $D_5\times\{0\}$ together with
\[
 \left(\frac{\sigma_1}2,\ldots,\frac{\sigma_5}2,
               \frac{\varepsilon\sqrt3}2\right),
 \qquad \prod_{i=1}^5\sigma_i=\varepsilon,
 \quad \varepsilon\in\{-1,1\}.
\]
There are $40+32=72$ vectors. Two new vectors in the same last-coordinate
layer differ in at least two of the first five signs, giving inner product
at most $1$. In opposite layers their first-five Hamming distance is odd,
and the inner product is $(1-d_H)/2\leq0$. A new vector has inner product
at most $1$ with every central root. Division by $\sqrt2$ gives the
claimed code. This is the familiar $E_6$ construction in explicit
coordinates, also discussed in \sref{110}{2}, Section~3.

Thus saturation plus an exact local rigidity certificate is insufficient
\emph{even in a neighboring dimension with the same root-system method}.
Nothing in the tangent argument excludes a disconnected component of the
configuration space containing more points. Nor do the computations
cover all real configurations of cardinality $41$ in dimension five.

\paragraph{Outcome.}
\textbf{Outcome: Exploratory approach with unresolved gap.}

The attempt obtains an exact covering-radius formula, a dimension-eight
hole-code completion and matching polynomial bound, and independently
checkable rigidity certificates. It also supplies an explicit
counterexample to the proposed local-to-global optimality principle.
These are restricted deductions or rederivations of known facts; no
novelty or improvement of a kissing-number record is claimed.

The full all-dimensional target remains open within this notebook.
Already in dimension five, a global exclusion of $41$ points or an
explicit larger configuration is missing. In higher dimensions, the root
and hole template is not an exhaustive parametrization and cannot be
promoted to an answer for every $n>4$.
% END RESEARCH ADDITION 110
\subsection{Sources}
\begin{itemize}\raggedright
\item[\textnormal{[S1]}]\hypertarget{source:110:1}{} Musin, The kissing number in four dimensions.
\url{https://arxiv.org/abs/math/0309430}.
{\small\textit{Catalogue formulation source.}}
\item[\textnormal{[S2]}]\hypertarget{source:110:2}{} Status source for Exact Kissing Numbers in Higher Dimensions.
\url{https://arxiv.org/html/2412.00937v3}.
{\small\textit{Catalogue research/status reference; not a proof endorsement.}}
% BEGIN NEW REFERENCES 110
\item[\textnormal{[E1]}]\hypertarget{extra:110:1}{}
Henry Cohn, Yang Jiao, Abhinav Kumar, and Salvatore Torquato,
``Rigidity of spherical codes,'' \emph{Geometry \& Topology}
\textbf{15} (2011), 2235--2273, arXiv:1102.5060v2.
Section~2, Theorem~2.2; Section~3, Corollary~3.4.
\url{https://arxiv.org/abs/1102.5060}.
\item[\textnormal{[E2]}]\hypertarget{extra:110:2}{}
Henry Cohn and Anqi Li, ``Improved kissing numbers in seventeen through
twenty-one dimensions,'' arXiv:2411.04916v3, 1 September 2026.
Theorem~1.1 and the subsequent note on dimension~19.
\url{https://arxiv.org/abs/2411.04916}.
\item[\textnormal{[E3]}]\hypertarget{extra:110:3}{}
Rustem Takhanov and Stanislav Yun, ``New lower bounds for kissing numbers
in dimensions $25$--$29$ and $31$,'' arXiv:2609.21591,
18 September 2026. Abstract consulted for recent claimed constructions;
not a proof dependency or independent certification.
\url{https://arxiv.org/abs/2609.21591}.
% END NEW REFERENCES 110
\end{itemize}

\end{document}
